\section{Introduction}

Much of extremal combinatorics is concerned with determining which global properties force the appearance of certain local substructures. One of the most important questions of this kind is the \emph{Tur\'an problem}, which asks how many edges an $n$-vertex graph (or hypergraph) can have without containing a given graph (or hypergraph) as a subgraph. The subject of this paper is another very central and closely related problem. A \emph{$(v,e)$-configuration} is a hypergraph with at most $v$ vertices and at least $e$ edges. Estimating the maximum number of edges an $n$-vertex $r$-uniform hypergraph can have without containing a $(v,e)$-configuration is an important topic in extremal graph theory. For example, when $e=\binom{v}{r}$, the problem becomes equivalent to determining the maximum possible number of edges without containing $K_v^{(r)}$, which is a notoriously difficult open problem for $r>2$.

Let us write $f(n, v, e)$ for the maximum number of edges in a $3$-uniform hypergraph on $n$ vertices which does not contain a $(v, e)$-configuration. In 1973, Brown, Erd\H os and S\' os \cite{BES73,BES73_2} initiated the study of this function for various values of $v$ and $e$. The following conjecture, named after them, is one of the most famous open problems in extremal combinatorics.

\begin{conjecture}[Brown--Erd\H os--S\'os conjecture~\cite{BES73,BES73_2}] \label{conj:BESConj}
    For every $e\geq 3$, we have $f(n,e+3,e)=o(n^2)$.
\end{conjecture}

The first result in this direction was the famous $(6, 3)$-theorem shown by Ruzsa and Szemer\'edi \cite{RS76}, who proved that $f(n, 6, 3)=o(n^2)$, thus establishing Conjecture~\ref{conj:BESConj} in the case $e=3$. 
This result has important implications in number theory. For instance, it implies Roth's theorem which states that any dense subset of first $n$ integers contains a three term arithmetic progression.
To this day, $e=3$ is the only instance of the Brown-Erd\H{o}s-S\'os conjecture which has been resolved. 

Therefore, a large amount of work has been put into proving approximate versions of Conjecture~\ref{conj:BESConj}. For instance, a classical result of S\'ark\"ozy and Selkow \cite{SS04} from 2004 is the following.
\begin{theorem}[S\'ark\"ozy and Selkow~\cite{SS04}]
For every $e \ge 3$, we have  $f(n, e+\lfloor \log_2 e\rfloor +2, e)=o(n^2)$.
\end{theorem}
For many years this result has been the state-of-the-art before Solymosi and Solymosi \cite{SS17} showed that $f(n, 14, 10)=o(n^2)$, thus improving the S\'ark\"ozy-Selkow theorem for $e=10$ which only gave $f(n, 15, 10)=o(n^2)$. In a recent breakthrough, Conlon, Gishboliner, Levanzov and Shapira~\cite{CGLS23} extended the ideas of Solymosi and Solymosi significantly to show that $f(n, e+O(\log e/\log\log e), e)=o(n^2)$. 

An important feature of all of the above proofs is their heavy use of the regularity lemma, both in the graph and the hypergraph setting. This means that the bounds they obtain on $f(n, v, e)$ are barely below quadratic. 

Motivated by various applications, Gowers and Long~\cite{GL21} asked to obtain a power-type improvement for the Brown--Erd\H os--S\'os problem.
They had made the surprising conjecture that $f(n, e+4, e)=O(n^{2-\eps})$ for all $e\geq 3$ and some $\eps>0$ potentially depending on $e$. 
They showed that even proving the special case $f(n,9,5)=O(n^{2-\eps})$ of this conjecture would answer a question of Ruzsa~\cite{ruzsa1993solving} on sets of integers avoiding solutions to a certain linear equation.
One cannot expect to strengthen Conjecture~\ref{conj:BESConj} to state $f(n, e+3, e)=O(n^{2-\eps})$ since there are $n$-vertex $3$-uniform hypergraphs with $n^{2-o(1)}$ edges avoiding $(e+3, e)$-configurations, for $e\in \{3, 4, 5, 7, 8\}$. More precisely, Ruzsa and Szemer\'edi~\cite{RS76} constructed a graph with $n$ vertices and $n^{2-o(1)}$ edges avoiding $(6, 3)$-configurations. Moreover, it is not hard to see that every $(7, 4)$- and $(8, 5)$-configuration contains a $(6, 3)$-configuration as a subhypergraph, and hence the Ruzsa-Szemer\'edi construction also shows that $f(n, 7, 4)\geq n^{2-o(1)}$ and $f(n, 8, 5)\geq n^{2-o(1)}$. Recently, lower bounds of the same type for $f(n, 10, 7)$ and $f(n, 11, 8)$ were also found by Ge and Shangguan~\cite{GS21}. Hence, in some sense, the conjecture of Gowers and Long is the strongest possible.

Conlon, Gishboliner, Levanzov and Shapira~\cite{CGLS23}, and later, Shapira and Tyomkyn~\cite{ST23} reiterated the following approximate version of the Gowers-Long conjecture. 
What is the smallest function $d=d(e)$ for which $f(n, e+d, e)=O(n^{2-\eps})$ for some $\eps = \eps(e)>0$? In this paper we prove the following, which can be thought of as a version of the S\'ark\"ozy-Selkow theorem with power saving.

\begin{theorem}\label{thm:power saving}
For every $e \ge 3$, there exists some $\eps>0$ such that $f(n, e+\lfloor \log_2 e\rfloor+38, e)=O(n^{2-\eps})$.
\end{theorem}

We remark that the weaker bound $f(n,e+O(\log_2 e),e)=O(n^{2-\eps})$ was independently proved by different groups of authors. Indeed, a bound of this form was obtained by Conlon~\cite{Conprivate}, by Gishboliner, Levanzov and Shapira \cite{GLSprivate} and by Gao et. al. \cite{Gaoetal}, using different techniques from one another. The purpose of our paper is to show that we obtain power saving already near the S\'ark\"ozy-Selkow bound, at the cost of a small additive constant. Note that in order to obtain power saving in the S\'ark\"ozy-Selkow bound when $e = 3$, an additive constant is necessary because, as remarked earlier, there exist $n$-vertex $3$-uniform hypergraphs with $n^{2-o(1)}$ edges avoiding a $(6, 3)$-configuration.

The value of our small additive constant comes from the fact that while our methods allow us to construct large configurations whose number of edges is not much smaller than the number of vertices, they do not allow us to control the exact number of edges easily. Hence, we have to perform a final cleaning step in which we take edges out of our configuration until exactly $e$ edges remain, and this is where the constant $38$ comes from. 

\subsection{Proof overview and notation}

A natural way to think about the Brown-Erd\H{o}s-S\'os conjecture is in terms of the deficiency. Namely, for a 3-uniform hypergraph $F$, we define the \textit{deficiency} of $F$, denoted by $\Delta(F)$, as the difference between the number of vertices and edges of $F$, i.e. $\Delta(F)=v(F)-e(F)$. Throughout the paper, we will often consider deficiencies of subgraphs of $F$. So for a set of vertices $U\subset V(F)$, we define the deficiency of $U$ as the deficiency of the subgraph of $F$ induced by $U$ and write $\Delta(U)=\Delta(F[U])$. In this language, Conjecture~\ref{conj:BESConj} asks to show that in a 3-uniform hypergraph $\cH$ with $\Omega(n^2)$ edges one can always find a subhypergraph $F\subseteq \cH$ with $e(F)=e$ and $\Delta(F)=3$.
